% Auto-generated from ch02-advanced.mdx — edit ch02-advanced.mdx then re-run mdx2tex.mjs
\chapter{第 2 章 参数化装饰器与类装饰器}
\label{ch:ch02-advanced}

\begin{tcolorbox}[colback=blue!5,colframe=blue!40,title={📘 本章地图}]
第 1 章的装饰器只能"裸用"（\textbackslash\{\}textbackslash\textbackslash\{\}\{\textbackslash\{\}\}texttt\textbackslash\{\}\{@log\textbackslash\{\}\}）。本章学会带参数的装饰器、装饰类的装饰器，以及多个装饰器堆叠时的执行顺序。

\begin{itemize}
\item 2.1 带参装饰器：为什么多一层函数 - 2.2 Recipe：@retry(times=N) - 2.3 类装饰器：装饰整个类 - 2.4 装饰器堆叠：自下而上应用 - 2.5 常见陷阱 - 2.6 要点总结
\end{itemize}
\end{tcolorbox}
\section{2.1 带参装饰器：为什么多一层函数}
第 1 章的 \textbackslash\{\}texttt\{@log\} 写死了行为。如果想 \textbackslash\{\}texttt\{@retry(times=3)\} 这样\textbackslash\{\}textbf\{带参数\}，必须多包一层函数：


\begin{codeblock}{python}
# 不带参：@deco           → deco(func) 返回 wrapper
# 带参：  @deco(times=3)  → deco(times=3)(func) 返回 wrapper
\end{codeblock}

也就是说，带参装饰器是\textbackslash\{\}textbf\{返回真正装饰器的工厂\}。通用模板：


\begin{codeblock}{python}
def deco_factory(params):
    """外层：接收装饰器参数，返回真正的装饰器。"""
    def real_decorator(func):
        """中层：真正的装饰器，输入 func，输出 wrapper。"""
        @functools.wraps(func)
        def wrapper(*args, **kwargs):
            """内层：被装饰后的函数。"""
            ...  # 可以使用 params、func
        return wrapper
    return real_decorator
\end{codeblock}

函数嵌套三层，理解关键是：\textbackslash\{\}textbf\{每次 \textbackslash\{\}texttt\{@deco(...)\} 先执行 \textbackslash\{\}texttt\{deco(...)\}，把结果当装饰器用。\}

\section{2.2 Recipe：@retry(times=N)}
<div className="callout recipe">

\subsubsection{🍳 Recipe 2-1：失败自动重试}
\textbackslash\{\}textbf\{Problem\}：调用不稳定的外部 API 时，想自动重试 N 次再放弃。

\textbackslash\{\}textbf\{Solution\}：


\begin{codeblock}{python}
# file: recipes/retry.py
import time
import functools

def retry(times: int = 3, delay: float = 0.5):
    """失败重试装饰器。

    Args:
        times: 最大重试次数（含首次调用）。
        delay: 每次重试前等待的秒数，避免雪崩。
    """
    def decorator(func):
        @functools.wraps(func)
        def wrapper(*args, **kwargs):
            last_exc = None
            for attempt in range(1, times + 1):
                try:
                    return func(*args, **kwargs)
                except Exception as exc:
                    last_exc = exc
                    print(f"[retry] 第 {attempt}/{times} 次失败：{exc}")
                    if attempt < times:
                        time.sleep(delay)  # 退避等待
            # 全部失败，向外抛出最后一次异常
            raise last_exc
        return wrapper
    return decorator

import random

@retry(times=3, delay=0.1)
def flaky_api() -> str:
    if random.random() < 0.7:
        raise ConnectionError("网络抖动")
    return "ok"

print(flaky_api())
\end{codeblock}

\textbackslash\{\}textbf\{Discussion\}：

\begin{itemize}
\item 三层嵌套对应 \textbackslash\{\}texttt\{retry(times=3)(func)(args)\}。
\item 用 \textbackslash\{\}texttt\{time.sleep(delay)\} 做简单退避；生产环境建议\textbackslash\{\}textbf\{指数退避\}（exponential backoff）。
\item 最后一次重试失败才 \textbackslash\{\}texttt\{raise\}，中间失败只打印日志。
\end{itemize}
</div>

\section{2.3 类装饰器：装饰整个类}
装饰器不仅能装饰函数，也能装饰\textbackslash\{\}textbf\{类\}。最常见的用途：

\begin{itemize}
\item 给类自动注册到某个全局表（如 ORM 模型、插件系统）
\item 把类改造为单例（singleton）
\item 自动给类加方法/属性
\end{itemize}

\begin{codeblock}{python}
# file: examples/singleton.py
import functools

def singleton(cls):
    """把类改造成单例：cls() 永远返回同一个实例。"""
    _instances = {}

    @functools.wraps(cls, updated=())
    def wrapper(*args, **kwargs):
        if cls not in _instances:
            _instances[cls] = cls(*args, **kwargs)
        return _instances[cls]

    return wrapper

@singleton
class AppConfig:
    def __init__(self):
        self.env = "production"

a = AppConfig()
b = AppConfig()
assert a is b  # True，同一个实例
\end{codeblock}

图 2-1 展示了装饰器相关类之间的继承与组合关系：

\textbackslash\{\}begin\{figure\}[htbp]\textbackslash\{\}centering\textbackslash\{\}includegraphics[width=0.85\textbackslash\{\}textwidth]\{figures/decorator-class.svg\}\textbackslash\{\}caption\{图 2-1：装饰器类层次\}\textbackslash\{\}end\{figure\}

注意：类装饰器接收的参数是 \textbackslash\{\}textbf\{类对象本身\}（而不是类的某个实例）。

\section{2.4 装饰器堆叠：自下而上应用}
多个装饰器可以堆叠在同一个函数上：


\begin{codeblock}{python}
@timer
@retry(times=3)
@log
def fetch(url): ...
\end{codeblock}

\textbackslash\{\}textbf\{执行顺序\}规则只有一句话：\textbackslash\{\}textbf\{离函数最近的装饰器最先应用（自下而上）\}。

上面的例子等价于：


\begin{codeblock}{python}
fetch = timer(retry(times=3)(log(fetch)))
\end{codeblock}

应用顺序是 \textbackslash\{\}texttt\{log → retry → timer\}；而\textbackslash\{\}textbf\{执行顺序\}（调用 \textbackslash\{\}texttt\{fetch()\} 时 wrapper 的堆叠顺序）正好相反：\textbackslash\{\}texttt\{timer → retry → log → 原 fetch\}。

记忆法：\textbackslash\{\}textbf\{装饰像洋葱，外到内层层包裹\}。

\subsection{装饰器类型速查}
\begin{table}[htbp]
\centering
\begin{tabularx}{\textwidth}{X X X X}
\toprule
类型 & 语法 & 层数 & 典型场景 \\
\midrule
无参函数装饰器 & \textbackslash\{\}texttt\{@deco\} & 2 层 & \textbackslash\{\}texttt\{@log\}, \textbackslash\{\}texttt\{@timer\} \\
带参函数装饰器 & \textbackslash\{\}texttt\{@deco(x=1)\} & 3 层 & \textbackslash\{\}texttt\{@retry(times=3)\}, \textbackslash\{\}texttt\{@lru\_cache(maxsize=64)\} \\
类装饰器 & \textbackslash\{\}texttt\{@deco\} 在 \textbackslash\{\}texttt\{class\} 上 & 2 层 & \textbackslash\{\}texttt\{@singleton\}, \textbackslash\{\}texttt\{@dataclass\} \\
堆叠装饰器 & 多个 \textbackslash\{\}texttt\{@\} & 多套 & 组合日志+重试+计时 \\
\bottomrule
\end{tabularx}
\end{table}

\section{2.5 常见陷阱}
<div className="callout pitfall">

\subsubsection{⚠️ 陷阱 2：装饰器堆叠顺序写反导致行为错乱}

\begin{codeblock}{python}
@log        # 外层
@retry(3)   # 内层
def api(): ...
\end{codeblock}

如果你期望"每次重试都打日志"，这个顺序是对的（外层 log 包裹整个 retry 循环，只会打一次）。如果你期望"每次尝试都打日志"，必须把 \textbackslash\{\}texttt\{@log\} 写在 \textbackslash\{\}texttt\{@retry\} \textbackslash\{\}textbf\{里面\}：


\begin{codeblock}{python}
@retry(3)
@log
def api(): ...
\end{codeblock}

\textbackslash\{\}textbf\{规则\}：想让某个 wrapper 对"每次重试"生效，就让它更靠近被装饰函数。

\subsubsection{⚠️ 陷阱 3：类装饰器忘记保留类元信息}
装饰类时 \textbackslash\{\}texttt\{functools.wraps(cls)\} 默认不会复制类的 \textbackslash\{\}texttt\{\_\_doc\_\_\}/\textbackslash\{\}texttt\{\_\_name\_\_\}，需要加 \textbackslash\{\}texttt\{updated=()\} 参数：


\begin{codeblock}{python}
@functools.wraps(cls, updated=())
def wrapper(*args, **kwargs): ...
\end{codeblock}

</div>

\section{2.6 要点总结}
\begin{tcolorbox}[colback=green!5,colframe=green!40,title={✅ 核心要点}]
\begin{itemize}
\item 带参装饰器 = "返回装饰器的工厂"，\textbackslash\{\}textbackslash\textbackslash\{\}\{\textbackslash\{\}\}textbf\textbackslash\{\}\{三层嵌套\textbackslash\{\}\}：工厂 → 装饰器 → wrapper。 - 类装饰器接收\textbackslash\{\}textbackslash\textbackslash\{\}\{\textbackslash\{\}\}textbf\textbackslash\{\}\{类对象\textbackslash\{\}\}，常见用途是单例、注册、改造类。 - 多个装饰器堆叠：\textbackslash\{\}textbackslash\textbackslash\{\}\{\textbackslash\{\}\}textbf\textbackslash\{\}\{自下而上\textbackslash\{\}\}应用（离函数最近最先执行装饰），运行时\textbackslash\{\}textbackslash\textbackslash\{\}\{\textbackslash\{\}\}textbf\textbackslash\{\}\{外→内\textbackslash\{\}\}包裹。 - 用 \textbackslash\{\}textbackslash\textbackslash\{\}\{\textbackslash\{\}\}texttt\textbackslash\{\}\{functools.wraps\textbackslash\{\}\} 保留元信息；装饰类时加 \textbackslash\{\}textbackslash\textbackslash\{\}\{\textbackslash\{\}\}texttt\textbackslash\{\}\{updated=()\textbackslash\{\}\}。
\end{itemize}
\end{tcolorbox}
行内公式示例：装饰器的数学抽象是 \$T: f \textbackslash\{\}mapsto f'\$，其中 \$f'\$ 是 wrapper。

块级公式示例：带参装饰器的三次调用可以写成：

\textbackslash\{\}[ D(a, b)(f)(x) = W\_\{a,b\}\textbackslash\{\}bigl(f(x)\textbackslash\{\}bigr) \textbackslash\{\}]

\section{2.7 延伸阅读}
\begin{itemize}
\item 第 1 章 · 装饰器基础
\item 附录 A · functools 速查
\item Python 官方文档 \textbackslash\{\}texttt\{functools\}：https://docs.python.org/3/library/functools.html
\end{itemize}
---

\textbackslash\{\}textbf\{本章参考文献\}

\begin{enumerate}
\item Python Software Foundation. "functools — Higher-order functions and operations on callable objects." https://docs.python.org/3/library/functools.html
\end{enumerate}
